\AImatter

\begin{ReportAiUse}{6}

\bibitem{AI1}
GitHub Copilot (Claude Opus 4.5, 2026年2月版本)\\
\textbf{查询1:} 我们提出了使用贝叶斯方法估计DWTS粉丝投票数的思路。由于数据集中只有淘汰结果而没有具体投票数，需要从淘汰模式反推投票分布。请帮助实现这个贝叶斯投票估计模型。\\
\textbf{输出:} AI根据我们的思路实现了层次贝叶斯模型代码框架：
\begin{lstlisting}[language=python]
class BayesianVoteModel:
    def __init__(self, weekly_data):
        self.weekly_data = weekly_data
        # Prior: votes follow log-normal distribution
        self.vote_prior_mu = 50000
        self.vote_prior_sigma = 30000
        
    def estimate_votes(self, week_data):
        # Posterior estimation using elimination signals
        # P(votes | elimination) proportional to 
        # P(elimination | votes) * P(votes)
        scores = week_data['total_score'].values
        eliminated = week_data['is_eliminated'].values
        
        # Use score-vote correlation as likelihood
        vote_estimates = self._bayesian_update(scores, eliminated)
        return vote_estimates
\end{lstlisting}
AI实现的这个代码框架成功地将我们的理论思路转化为可执行的模型，后续我们对参数进行了调优和验证。

\bibitem{AI2}
GitHub Copilot (Claude Opus 4.5, 2026年2月版本)\\
\textbf{查询2:} 我们设计了使用随机森林分析名人特征（年龄、行业、搭档经验）对比赛表现影响的方案。请实现特征重要性分析的代码，包括模型训练和结果排序。\\
\textbf{输出:} AI实现了我们设计的特征重要性分析框架：
\begin{lstlisting}[language=python]
class RandomForestAnalyzer:
    def __init__(self, data):
        self.model = RandomForestRegressor(
            n_estimators=100, max_depth=10, random_state=42)
        
    def get_feature_importance(self, X, y, feature_names):
        self.model.fit(X, y)
        importance = pd.DataFrame({
            'feature': feature_names,
            'importance': self.model.feature_importances_
        }).sort_values('importance', ascending=False)
        return importance
\end{lstlisting}
基于AI实现的代码，我们成功分析出\texttt{week}（0.48）、\texttt{age}（0.18）和\texttt{partner\_experience}（0.13）是评委打分的关键因素。

\bibitem{AI3}
GitHub Copilot (Claude Opus 4.5, 2026年2月版本)\\
\textbf{查询3:} 我们设计了一个动态加权投票系统（戏剧弧线系统），核心思路是：早期偏重评委分（52\%），中期平衡（50-50），后期偏重观众票（48\%）。请实现这个三阶段动态权重系统的代码。\\
\textbf{输出:} AI根据我们的设计方案实现了系统代码：
\begin{lstlisting}[language=python]
class DramaticArcSystem(VotingSystem):
    def get_weights(self, week_num, total_weeks):
        progress = week_num / total_weeks
        if progress <= 0.3:  # Early stage: Judge-focused
            judge_weight = 0.52
        elif progress <= 0.6:  # Mid stage: Balanced
            judge_weight = 0.50
        else:  # Late stage: Fan-focused
            judge_weight = 0.48
        return judge_weight, 1 - judge_weight
\end{lstlisting}
AI实现的代码成功体现了我们的设计理念，经过后续调优，该系统达到了86\%的公平性和11.1\%的争议率。

\bibitem{AI4}
GitHub Copilot (Claude Opus 4.5, 2026年2月版本)\\
\textbf{查询4:} 我们设计了综合敏感性分析方案，包括参数扰动分析和蒙特卡罗模拟。请实现这个分析框架的代码，能够系统性地测试各参数对模型结果的影响。\\
\textbf{输出:} AI实现了我们设计的多方法敏感性分析框架：
\begin{lstlisting}[language=python]
class SensitivityAnalyzer:
    def parameter_sensitivity(self, param_name, values):
        results = []
        for val in values:
            self.config[param_name] = val
            metrics = self.evaluate_system()
            results.append({
                'param_value': val,
                'fairness': metrics['fairness'],
                'excitement': metrics['excitement']
            })
        return pd.DataFrame(results)
    
    def monte_carlo_analysis(self, n_simulations=1000):
        # Randomly perturb all parameters within +/-20%
        outcomes = []
        for _ in range(n_simulations):
            perturbed = self._perturb_params()
            outcome = self.evaluate_with_params(perturbed)
            outcomes.append(outcome)
        return self._compute_statistics(outcomes)
\end{lstlisting}

\bibitem{AI5}
GitHub Copilot (Claude Opus 4.5, 2026年2月版本)\\
\textbf{查询5:} 需要创建高质量的图表展示我们的分析结果。请实现可视化函数，包括：系统对比的条形图（4个子图展示Fairness、Excitement、Consistency、Simplicity）、权重变化曲线图、敏感性分析的龙卷风图等。图表需要publication-quality，DPI 300。\\
\textbf{输出:} AI实现了我们需要的可视化代码：
\begin{lstlisting}[language=python]
def plot_system_comparison(comparison_df, save_path):
    fig, axes = plt.subplots(2, 2, figsize=(14, 10))
    metrics = ['Fairness', 'Excitement', 'Consistency', 'Simplicity']
    
    for ax, metric in zip(axes.flat, metrics):
        data = comparison_df.sort_values(metric, ascending=True)
        colors = [COLORS.get(s, '#888888') for s in data['System']]
        ax.barh(data['System'], data[metric], color=colors)
        ax.set_xlabel(metric)
        ax.set_xlim(0, 1)
    
    plt.tight_layout()
    plt.savefig(save_path, dpi=300, bbox_inches='tight')
\end{lstlisting}
AI生成的可视化代码产生了符合论文要求的高质量图表，清晰展示了我们的研究结果。

\bibitem{AI6}
GitHub Copilot (Claude Opus 4.5, 2026年2月版本)\\
\textbf{查询6:} 代码运行时出现ValueError: The truth value of a Series is ambiguous。错误出现在争议奖金计算部分，使用了min()函数比较Series。请帮助定位并修复这个bug。\\
\textbf{输出:} AI快速识别了问题并提供了修复方案：
\begin{lstlisting}[language=python]
# 原始有bug的代码:
df.loc[df['is_controversial'], 'bonus'] = min(
    df['score_gap'] * 0.15, 0.08)  # 错误: 比较Series

# 修复后的代码:
controversial_mask = df['is_controversial']
df.loc[controversial_mask, 'bonus'] = np.minimum(
    df.loc[controversial_mask, 'score_gap'] * 0.15, 0.08)
\end{lstlisting}
AI的修复解决了我们的问题，改用\texttt{np.minimum()}进行逐元素比较，使代码能够正确处理pandas Series对象。

\end{ReportAiUse}